% !TEX TS-program = pdflatex
% !TeX program = pdflatex
% !TEX encoding = UTF-8
% !TEX spellcheck = fr

\documentclass[xcolor=table]{beamer}


%\usepackage{fullpage}
%\usepackage[left=2.8cm,right=2.2cm,top=2 cm,bottom=2 cm]{geometry}
\setbeamersize{text margin left=10pt,text margin right=10pt}
\usepackage{amsmath,amssymb} 
\usepackage[T1]{fontenc}
\usepackage[utf8]{inputenc}
\usepackage[english,french]{babel}
\usepackage{txfonts}
\usepackage[]{graphicx}
\usepackage{multirow}
\usepackage{hyperref}
\usepackage{colortbl}
\usepackage{listings}
\usepackage{wrapfig}
\usepackage{multicol}
\usepackage[export]{adjustbox} %for images in table, also for frame

\hypersetup{
	colorlinks,
	urlcolor = blue
}

%\renewcommand{\baselinestretch}{1.5}

\def\supit#1{\raisebox{0.8ex}{\small\it #1}\hspace{0.05em}}

\AtBeginSection{%
	\begin{frame}
		\sectionpage
	\end{frame}
}

\newcommand{\rottext}[2]{%
	\rotatebox{90}{%
	\begin{minipage}{#1}%
		\raggedleft#2%
	\end{minipage}%
	}%
}

\usepackage{longtable}
\usepackage{tabu}


\institute{ %
École  nationale Supérieure d'Informatique (ESI, ex. INI), Algérie
}
\author[ \textbf{\footnotesize  \insertframenumber/\inserttotalframenumber} \hspace*{\fill} ESI (2019-2020)] %
{ARIES Abdelkrime}
%\titlegraphic{\includegraphics[height=1cm]{../img/esi-logo.png}%\hspace*{4.75cm}~


\date{Année unniversitaire: 2019/2020} %\today

\usetheme{Warsaw} % Antibes Boadilla Warsaw

\beamertemplatenavigationsymbolsempty

%\setbeamertemplate{headline}{}

\definecolor{lightblue}{HTML}{D0D2FF}
\definecolor{lightyellow}{HTML}{FFFFAA}
\definecolor{darkblue}{HTML}{0000BB}
\definecolor{olivegreen}{HTML}{006600}
\definecolor{violet}{HTML}{6600CC}

\newcommand{\keyword}[1]{\textcolor{red}{\bfseries\itshape #1}}
\newcommand{\expword}[1]{\textcolor{olivegreen}{#1}}
\newcommand{\optword}[1]{\textcolor{violet}{\bfseries #1}}

\makeatletter
\newcommand\mysphere{%
	\parbox[t]{10pt}{\raisebox{0.2pt}{\beamer@usesphere{item projected}{bigsphere}}}}
\makeatother

%\let\oldtabular\tabular
%\let\endoldtabular\endtabular
%\renewenvironment{tabular}{\rowcolors{2}{white}{lightblue}\oldtabular\rowcolor{blue}}{\endoldtabular}


\NoAutoSpacing %french autospacing after ":"

\def\graphpath{}

\newcommand{\changegraphpath}[1]{\def\graphpath{#1}}

\newcommand{\graphintable}[2]{%
	\includegraphics[height=#1,valign=t]{\graphpath #2}%
}

\newcommand{\vgraphpage}[2][.85\textheight]{%
	\begin{center}%
		\includegraphics[height=#1]{\graphpath #2}%
	\end{center}%
}

\newcommand{\hgraphpage}[2][\textwidth]{%
%	\begin{center}%
		\includegraphics[width=#1]{\graphpath #2}%
%	\end{center}%
}

\newcommand{\insertbibliography}[2]{
	\appendix
	\section*{Bibliographie}
	\nocite{#2}
	
	\begin{multicols*}{2}[\frametitle{\insertsection} \usebeamertemplate{frametitle}]%\usebeamertemplate*{frametitle}\frametitle{Références}
		\tiny
		\bibliography{#1}
		\bibliographystyle{acm}
	\end{multicols*}
}



\title[BWEB : 07- CSS] %
{Bureautique et Web \\Chapitre 07 : Développement Web\\ \slshape\small  CSS}  

\changegraphpath{../img/dev-web/}

\begin{document}
	
\begin{frame}
\frametitle{CSS}
\framesubtitle{Introduction}

\begin{itemize}
	\item CSS : Cascading Style Sheets
	\item un langage de feuille de style 
	\item il sert à décrire la présentation d'une page web
\end{itemize}

\end{frame}

\begin{frame}
\frametitle{CSS}
\framesubtitle{Fonctionnement}

\begin{figure}
	\centering
	\hgraphpage{rendering.pdf}
	\caption{Application de CSS \cite{mdn}}
\end{figure}
\begin{itemize}
	\item \keyword{DOM} : Document Object Model
\end{itemize}

\end{frame}


\begin{frame}
\frametitle{CSS}
\framesubtitle{Plan}

\begin{multicols}{2}
%	\small
\tableofcontents
\end{multicols}
\end{frame}

%===================================================================================
\section{Règles et cascade}
%===================================================================================

\subsection{Règles et valeurs}

\begin{frame}[fragile]
\frametitle{CSS : Règles et cascade}
\framesubtitle{Règles et valeurs : Règles}

\begin{minipage}{0.60\textwidth} 
\begin{itemize}
	\item Une règle se compose d'un sélecteur et d'une déclaration (propriété/valeur)
	\item \optword{Sélecteur} : permet de sélectionner les éléments sur lesquels appliquer le style souhaité
	\item \optword{Propriété} : différentes façons dont on peut mettre en forme un élément HTML
	\item \optword{Valeur} : permet de choisir une mise en forme parmi d'autres pour une propriété donnée
\end{itemize}
\end{minipage}
%
\begin{minipage}{0.38\textwidth}
\begin{exampleblock}{Structure d'une règle CSS}
\lstset{escapeinside=**}
\scriptsize\bfseries
\begin{lstlisting}[language={html}]
*sélecteur* {
  *proriété1*: valeur1;
  *proriété2*: valeur2;
}
\end{lstlisting}
\end{exampleblock}
\begin{figure}
	\centering
	\hgraphpage[.6\textwidth]{CSS-regle.pdf}
	\caption{Structure d'une règle CSS \cite{mdn}}
\end{figure}
\end{minipage}

\end{frame}

\begin{frame}[fragile]
\frametitle{CSS : Règles et cascade}
\framesubtitle{Règles et valeurs : Unités absolues}

\scriptsize\bfseries
\rowcolors{2}{lightblue}{lightyellow}
\begin{minipage}{0.50\textwidth} 
\begin{center}
	\begin{tabular}{lll}
		\rowcolor{darkblue}
		\color{white}Unité & \color{white}Nom & \color{white}Description \\
		cm & Centimètres & 1cm = 96px/2.54\\
		mm & Millimètres & 1mm = 1/10 cm\\
		Q & 1/4 millimètres & 1Q = 1/40 cm\\
		in & Pouces & 1in = 2.54cm = 96px\\
		pc & Picas & 1pc = 1/6 in\\
		pt & Points & 1pt = 1/72 in\\
		px & Pixels & 1px = 1/96 in\\
	\end{tabular}
\end{center}
\end{minipage}
%
\begin{minipage}{0.49\textwidth}
	\begin{center}
		\begin{tabular}{llp{.3\textwidth}}
			\rowcolor{darkblue}
			
			\color{white}Unité & \color{white}Nom & \textcolor{white}{Description} \\
			s & Secondes &  \\
			ms & Millisecondes & Un millième de seconde.\\
		\end{tabular}
	\end{center}
	\begin{center}
		\begin{tabular}{llp{.4\textwidth}}
			\rowcolor{darkblue}
			\color{white}Unité & \color{white}Nom & \color{white}Description \\
			Hz & Hertz & Nombre de fois par seconde.\\
			kHz & Kilohertz & 1000 Hertz.\\
		\end{tabular}
	\end{center}
\end{minipage}

%\scriptsize\bfseries
%\rowcolors{2}{lightblue}{lightyellow}
\begin{center}
	\begin{tabular}{lll}
		\rowcolor{darkblue}
		\color{white}Unité & \color{white}Nom & \color{white}Description \\
		deg & Degrés & Un cercle se divise en 360 degrés égaux. \\
		grad & Grades/Gradians & Un cercle se compose de 400 grades. \\
		rad & Radians & Un cercle se compose de $2\pi$ radians. \\
		turn & Tours & Un cercle se compose d'un tour. \\
	\end{tabular}
\end{center}

\end{frame}

\begin{frame}[fragile]
\frametitle{CSS : Règles et cascade}
\framesubtitle{Règles et valeurs : Unités relatives}

\scriptsize\bfseries
\rowcolors{2}{lightblue}{lightyellow}

\begin{center}
	\begin{tabular}{p{.1\textwidth}p{.8\textwidth}}
		\rowcolor{darkblue}
		\color{white}Unité & \color{white}Sens\\
		em & * Dans le cas de \keyword{font-size} : Taille de police du parent 
		
		* Dans d'autres cas comme \keyword{width} : Taille de police  de l'élément lui-même\\
		ex & La hauteur du caractère minuscule "x"\\
		ch & La largeur du caractère "0"\\
		rem & Taille de police de l'élément racine\\
		lh & Hauteur de ligne de l'élément\\
		vw & 1 \% de la largeur de la fenêtre\\
		vh & 1 \% de la longueur de la fenêtre\\
		vmin & 1 \% de la dimension minimale de la fenêtre\\
		vmax & 1 \% de la dimension maximale de la fenêtre\\
		\% & Pourcentage par rapport à l'élément parent\\
	\end{tabular}
\end{center}

\end{frame}

\subsection{Application de CSS}

\begin{frame}[fragile]
\frametitle{CSS : Règles et cascade}
\framesubtitle{Application de CSS}

\begin{itemize}
	\item Styles en ligne
	\begin{itemize}
		\item définir un style dans la balise de l'élément
		\item en utilisant l'attribut \keyword{style} de l'élément
	\end{itemize}
	\item Feuille de style interne
	\begin{itemize}
		\item définir une feuille de style dans une page HTML
		\item en utilisant la balise \keyword{<style>} dans \keyword{<head>}
	\end{itemize}
	\item Feuille de style externe
	\begin{itemize}
		\item définir une feuille de style séparée
		\item en utilisant la balise \keyword{<link>} dans \keyword{<head>}
	\end{itemize}
\end{itemize}

\end{frame}

\begin{frame}[fragile]
\frametitle{CSS : Règles et cascade}
\framesubtitle{Application de CSS : Style en ligne}

\begin{minipage}{0.60\textwidth} 
	\begin{itemize}
		\item En utilisant l'attribut \keyword{style} de la balise
		\item Mise en forme d'un seul élément
		\item Si on veut utiliser la même mise en forme pour plusieurs éléments, on doit la recopier
		\item Exemple, \expword{Si on veut utiliser la même couleur pour tous les éléments de type <p>, on doit recopier le style dans chaque éléments}
	\end{itemize}
\end{minipage}
%
\begin{minipage}{0.38\textwidth}
\begin{exampleblock}{Style en ligne}
\lstset{escapeinside=**}
\scriptsize\bfseries
\begin{lstlisting}[language={html}]
<html>
  <head>
    <title>en ligne</title>
  </head>
  <body>
    <p style="color:red;">
      Paragraphe en rouge
    </p>
  </body>
<html>
\end{lstlisting}
\end{exampleblock}
\end{minipage}
\end{frame}

\begin{frame}[fragile]
\frametitle{CSS : Règles et cascade}
\framesubtitle{Application de CSS : Feuille de style interne}

\begin{minipage}{0.60\textwidth} 
	\begin{itemize}
		\item En utilisant la balise \keyword{<style>} dans l'entête
		\item \keyword{type="text/css"} pour dire au navigateur qu'il s'agit de CSS
		\item Mise en forme des éléments d'une page HTML
		\item Si on veut utiliser la même feuille pour plusieurs pages, on doit la recopier
		\item Exemple, \expword{Si on veut utiliser la même couleur pour tous les éléments de type <p>, on doit recopier le style dans chaque éléments}
	\end{itemize}
\end{minipage}
%
\begin{minipage}{0.38\textwidth}
\begin{exampleblock}{Style en ligne}
\lstset{escapeinside=**}
\scriptsize\bfseries
\begin{lstlisting}[language={html}]
<html>
  <head>
    <title>interne</title>
    <style type="text/css">
      p {color: red;}
    </style>
  </head>
  <body>
    <p>Paragraphe1</p>
    <div>DIV1</div>
    <p>Paragraphe2</p>
    <p>Paragraphe3</p>
  </body>
<html>
\end{lstlisting}
\end{exampleblock}
\end{minipage}
\end{frame}

\begin{frame}[fragile]
\frametitle{CSS : Règles et cascade}
\framesubtitle{Application de CSS : Feuille de style externe}

\begin{minipage}{0.60\textwidth} 
	\begin{itemize}
		\item En utilisant la balise \keyword{<style>} dans l'entête
		\item \keyword{type="text/css"} peut être omis
		\item \keyword{rel="stylesheet"} pour dire au navigateur qu'il s'agit d'une feuille de style
		\item \keyword{href} l'emplacement de la feuille
		\item Mise en forme des éléments de plusieurs pages 
	\end{itemize}
\begin{minipage}{0.95\textwidth} 
\begin{exampleblock}{style.css}
	\lstset{escapeinside=**}
	\scriptsize\bfseries\vspace{-6pt}
\begin{lstlisting}[language={CSS}]
p {
  color: red;
}
\end{lstlisting}\vspace{-6pt}
\end{exampleblock}
\end{minipage}
\end{minipage}
%
\begin{minipage}{0.38\textwidth}
\begin{exampleblock}{Lier la page avec une feuille de style externe}
\lstset{escapeinside=**}
\scriptsize\bfseries
\begin{lstlisting}[language={html}]
<html>
  <head>
    <title>externe</title>
    <link rel="stylesheet" 
         href="style.css">
  </head>
  <body>
    <p>Paragraphe1</p>
    <div>DIV1</div>
    <p>Paragraphe2</p>
  </body>
<html>
\end{lstlisting}
\end{exampleblock}
\end{minipage}

\end{frame}

\subsection{Sélecteurs}

\begin{frame}
\frametitle{CSS : Règles et cascade}
\framesubtitle{Sélecteurs}

\begin{itemize}
	\item Les sélecteurs simples
	\begin{itemize}
		\item \optword{Sélecteur universel}
		\item \optword{Sélecteurs de type}
		\item \optword{Sélecteurs de classe}
		\item \optword{Sélecteurs d'identifiant}
		\item \optword{Sélecteurs d'attribut}
	\end{itemize}
	\item Les combinateurs
	\begin{itemize}
		\item \optword{Sélecteurs d'éléments descendants}
		\item \optword{Sélecteurs d'éléments fils}
		\item \optword{Sélecteurs de voisin}
		\item \optword{Sélecteurs de voisin direct}
	\end{itemize}
	\item \optword{Les pseudo-classes}
	\item \optword{Les pseudo-éléments}
\end{itemize}

\end{frame}

\begin{frame}[fragile]
\frametitle{CSS : Règles et cascade}
\framesubtitle{Sélecteurs : Sélecteurs simples}

\begin{minipage}{0.60\textwidth} 
	\begin{itemize}
		\item \optword{universel} sélectionne tous les éléments en utilisant le mot clé \keyword{*}
		\item \optword{de type} l'élément est ciblé par so type
		\item \optword{de classe} l'élément est ciblé par la valeur de son attribut \keyword{class} en précédant cette valeur par \keyword{.}
		\item \optword{d'identifiant} l'élément est ciblé par la valeur de son attribut \keyword{id} en précédant cette valeur par \keyword{\#}
		\item \optword{d'attribut} l'élément est ciblé par la valeur d'un attribut donné en utilisant la syntaxe \expword{[attr operateur valeur]}
	\end{itemize}
\end{minipage}
%
\begin{minipage}{0.38\textwidth}
\begin{exampleblock}{selecteur\_simple.html}
%\lstset{escapeinside=**}
\scriptsize\bfseries
\begin{lstlisting}[language={CSS}]
* {font-size: 12pt;}
p {color: magenta;}
.p {color: green;}
#p {color: blue;}
input[type=text] {
  color: red;
}
\end{lstlisting}
\end{exampleblock}
\end{minipage}
\end{frame}


\begin{frame}[fragile]
\frametitle{CSS : Règles et cascade}
\framesubtitle{Sélecteurs : Combinateurs}

\begin{minipage}{0.60\textwidth} 
	\begin{itemize}
		\item \optword{d'éléments descendants} on sépare l'élément et son descendant par un espace
		\item \optword{d'éléments fils} on sépare l'élément et son fils par \keyword{>}
		\item \optword{de voisin} on sépare l'élément et son voisin par \keyword{\textasciitilde}
		\item \optword{de voisin direct} on sépare l'élément et son voisin par \keyword{+}
	\end{itemize}
\end{minipage}
%
\begin{minipage}{0.38\textwidth}
\begin{exampleblock}{combinateur.html}
\lstset{escapeinside=**}
\scriptsize\bfseries
\begin{lstlisting}[language={CSS}]
#id1 p {color:blue;}
#id2 > p {color:red;}
h1 ~ p {color:green;}
h2 + p {color:orange;}
\end{lstlisting}
\end{exampleblock}
\end{minipage}
\end{frame}

\begin{frame}[fragile]{Code}
\frametitle{CSS : Règles et cascade}
\framesubtitle{Sélecteurs : Pseudo-classes}

\begin{minipage}{0.60\textwidth} 
	\begin{itemize}
		\item l'élément est ciblé par un état spécifique
		\item \keyword{:hover} la souris est positionnée sur l'élément
		\item \keyword{:visited} un lien visité
		\item \keyword{:link} un lien non visité
		\item \keyword{:active} un élément activé par l'utilisateur.
		\item \keyword{:checked} un bouton radio, case à couché ou une option sélectionnée
		\item \keyword{:disabled} un élément désactivé
		\item \keyword{:enabled} un élément activé
		\item \keyword{:required} un champs obligatoire
	\end{itemize}
\end{minipage}
%
\begin{minipage}{0.38\textwidth}
\begin{exampleblock}{pseudo-classes1.html}
%\lstset{escapeinside=**}
\scriptsize\bfseries
\begin{lstlisting}[language={CSS}]
a:link {color: blue;}
a:visited {color: purple;}
a:hover {color: yellow;}
a:active {color: red;}
input:checked {color:blue;}
input:disabled {color:grey;}
input:required {color:red;}
\end{lstlisting}
\end{exampleblock}
\end{minipage}

\end{frame}

\begin{frame}[fragile]
\frametitle{CSS : Règles et cascade}
\framesubtitle{Sélecteurs : Pseudo-classes (2)}

\begin{minipage}{0.60\textwidth} 
	\begin{itemize}
		\item \keyword{:not(s)} un élément qui n'est pas un élément sélectionné par le sélecteur \expword{s}
		\item \keyword{:nth-child(an+b)} un élément qui possède an+b-1 éléments voisins (au même niveau) avant lui 
		\item \keyword{:nth-of-type(an+b)} des éléments d'un type donné, en fonction de leur position au sein d'un groupe de frères et sœurs.
		\item \keyword{:out-of-range} un élément \keyword{<input>} lorsque sa valeur est en dehors de l'intervalle autorisé par les attributs \keyword{min} et \keyword{max}.
	\end{itemize}
\end{minipage}
%
\begin{minipage}{0.38\textwidth}
\begin{exampleblock}{pseudo-classes2.html}
%\lstset{escapeinside=**}
\scriptsize\bfseries\vspace{-6pt}
\begin{lstlisting}[language={html}]
*:not(h1){font-style: italic;}
tr:nth-child(2n) {
  background: lightblue;
}
tr:nth-child(2n+1) {
  background: lightgreen;
}
#id1 p:nth-child(2n+1) {
  color: red;
}
#id2 p:nth-of-type(2n+1) {
  color: blue;
}
input:out-of-range {
  background: red;
}
\end{lstlisting}\vspace{-6pt}
\end{exampleblock}
\end{minipage}
\end{frame}

\begin{frame}[fragile]
\frametitle{CSS : Règles et cascade}
\framesubtitle{Sélecteurs : Pseudo-éléments}

\begin{minipage}{0.60\textwidth} 
	\begin{itemize}
		\item Mettre en forme certaines parties de l'élément ciblé par la règle.
		\item \keyword{::first-line} la première ligne d'un élément
		\item \keyword{::first-letter} la première lettre d'un élément
		\item \keyword{::selection} la portion sélectionnée d'un élément
		\item \keyword{::before} insérer une chose avant le contenu d'un élément
		\item \keyword{::after} insérer une chose après le contenu d'un élément
	\end{itemize}
\end{minipage}
%
\begin{minipage}{0.38\textwidth}
\begin{exampleblock}{pseudo-element.html}
\lstset{escapeinside=**}
\scriptsize\bfseries
\begin{lstlisting}[language={html}]
p#id1::first-letter {
  color: blue; }
p#id2::first-line {
  color: red; }
p#id3::selection {
  color: green; }
p#id4::before {
  content: "avant ";
  color: orange;
}
p#id4::after {
  content: " *après*";
  color: magenta;
}
\end{lstlisting}
\end{exampleblock}
\end{minipage}
\end{frame}

\subsection{Cascade et héritage}

%\begin{frame}
%\frametitle{CSS : Les règles et la cascade}
%\framesubtitle{Cascade et héritage}
%
%\begin{itemize}
%	\item 
%\end{itemize}
%
%\end{frame}

\begin{frame}[fragile]
\frametitle{CSS : Règles et cascade}
\framesubtitle{Cascade et héritage : Héritage}

\begin{minipage}{0.60\textwidth}
\begin{itemize}
	\item Si un élément contient un autre élément, ce dernier hérite les propriétés e son parent
	\item Sauf si la propriété a été redéfinie
	\item Par exemple, \expword{Si on définit la couleur pour l'élément <body>, tous les éléments dedans vont hériter cette propriété}
\end{itemize}
\end{minipage}
%
\begin{minipage}{0.38\textwidth}
\begin{exampleblock}{Héritage dans CSS}
\lstset{escapeinside=**}
\scriptsize\bfseries
\begin{lstlisting}[language={HTML5}]
<html>
  <head>
    <title>interne</title>
    <style type="text/css">
      body {color: red;}
    </style>
  </head>
  <body>
    <p>Paragraphe1</p>
    <div>DIV1</div>
    <p>Paragraphe2</p>
    <p>Paragraphe3</p>
  </body>
<html>
\end{lstlisting}
\end{exampleblock}
\end{minipage}
\end{frame}

\begin{frame}[fragile]
\frametitle{CSS : Règles et cascade}
\framesubtitle{Cascade et héritage : Cascade}

\begin{minipage}{0.60\textwidth}
	\begin{itemize}
		\item L'origine de la feuille
		\begin{itemize}
			\item \optword{Agent utilisateur} : chaque navigateur fournit une feuille de style par défaut
			\item \optword{Auteur} : celui qui a créé le site web 
			\item \optword{Utilisateur} : celui qui utilise le site web 
		\end{itemize}
		\item L'algorithme de la cascade
		\begin{enumerate}
			\item Filtrer les règles des différentes sources pour garder celles applicables sur un élément
			\item Trier les règles en se basant sur leurs origines et importances.
			\item Appliquer les règles les plus importantes 
		\end{enumerate}
	\end{itemize}
\end{minipage}
%
\begin{minipage}{0.38\textwidth}
\begin{center}
	\scriptsize\bfseries
	\rowcolors{2}{lightblue}{lightyellow}
	\begin{tabular}{ll}
		\rowcolor{darkblue}
		\color{white}Origine & \color{white}Importance\\
		Agent utilisateur & normale\\
		Utilisateur & normale\\
		Auteur & normale\\
		Auteur & !important\\
		Utilisateur & !important\\
		Agent utilisateur & !important\\
	\end{tabular}
\end{center}
\begin{exampleblock}{Exemple d'une règle importante}
\lstset{escapeinside=**}
\tiny\bfseries
\begin{lstlisting}[language={HTML5}]
p {
  color: red;
  font-size: 14pt !important;
}
\end{lstlisting}
\end{exampleblock}
\end{minipage}
\end{frame}


\begin{frame}[fragile]
\frametitle{CSS : Règles et cascade}
\framesubtitle{Cascade et héritage : Spécificité}

\begin{minipage}{0.60\textwidth}
	\begin{itemize}
		\item On calcule la spécificité en se basant sur le sélecteur
		\item une règle en ligne : spécificité \optword{1000} 
		\item le nombre des IDs dans le sélecteur multiplié par \optword{100}
		\item le nombre des classes, les sélecteurs d'attributs et pseudo-classes dans le sélecteur multiplié par \optword{10}
		\item le nombre des éléments et pseudo-éléments dans le sélecteur
		\item on fait la somme
	\end{itemize}
\end{minipage}
%
\begin{minipage}{0.38\textwidth}
\begin{exampleblock}{Exemple de spécificité}
\lstset{escapeinside=!!}
\scriptsize\bfseries
\begin{lstlisting}[language={html}]
p a {}/*!spécificité! = 002*/
p .a {}/*!spécificité! = 011*/
p #a {}/*!spécificité! = 101*/
/*!spécificité! = 112*/
div#a p.c {}
/*!spécificité! = 011*/
input[type=text] {}
/*!spécificité! = 012*/
a:hover img {}
\end{lstlisting}
\end{exampleblock}
\end{minipage}
\end{frame}

%===================================================================================
\section{Mise en page}
%===================================================================================

\subsection{Modèle en boîte}

\begin{frame}[fragile]
\frametitle{CSS : Mise en page}
\framesubtitle{Modèle en boîte : Rendu et disposition (1)}

\begin{minipage}{0.60\textwidth}
	\begin{itemize}
		\item On définit le rendu et la disposition des éléments en utilisant la propriété \keyword{display}
		\item \optword{none} : l'élément est caché
		\item \optword{block} : l'élément est défini comme bloc
		\item \optword{inline} : l'élément est défini comme en ligne
		\item \optword{inline-block} : l'élément est défini comme en ligne, mais on peut définir la hauteur et la largeur
	\end{itemize}
\end{minipage}
%
\begin{minipage}{0.38\textwidth}
\begin{exampleblock}{display.html}
\lstset{escapeinside=**}
\scriptsize\bfseries
\begin{lstlisting}[language={CSS}]
p {
  background: red;
  height: 50px;
}
#p1 {display:none;}
#p2 {display:block;}
#p3 {display:inline;}
#p4 {display:inline-block;}
\end{lstlisting}
\end{exampleblock}
\end{minipage}
\end{frame}

\begin{frame}[fragile]
\frametitle{CSS : Mise en page}
\framesubtitle{Modèle en boîte : Rendu et disposition (2)}

\begin{minipage}{0.60\textwidth}
	\begin{itemize}
		\item \keyword{display: flex;}
		\item les éléments d'un conteneur flexible peuvent être disposés dans n'importe quelle direction
		\item l'attribut \keyword{flex-direction} définit la direction des éléments fils
		\begin{itemize}
			\item \optword{row} : la direction du texte
			\item \optword{row-inverse} : l'inverse de la direction du texte
			\item \optword{column} : de haut en bas
			\item \optword{column-inverse} : de bas en haut
		\end{itemize}
	\item \keyword{flex-wrap: wrap;} : les éléments peuvent être affichés sur plusieurs lignes 
	\end{itemize}
\end{minipage}
%
\begin{minipage}{0.38\textwidth}
\begin{exampleblock}{display.html}
\lstset{escapeinside=**}
\tiny\bfseries
\begin{lstlisting}[language={CSS}]
div {
  display: flex;
  flex-wrap: wrap;
}
.c {width: 30%;}
#div5 {
  flex-direction:row;}
#div6 {
  flex-direction:row-reverse;}
#div7 {
  flex-direction:column;}
#div8 {
  flex-direction:column-reverse;}
\end{lstlisting}
\end{exampleblock}
\end{minipage}
\end{frame}

\begin{frame}[fragile]
\frametitle{CSS : Mise en page}
\framesubtitle{Modèle en boîte : Composition d'une boîte}

\begin{minipage}{0.60\textwidth}
	\begin{itemize}
		\item \keyword{margin} : la marge entre l'élément et ses voisins
		\item \keyword{border} : la bordure de l'élément 
		\item \keyword{paddin} : la marge entre la bordure et le contenu
		\item \keyword{width} : la largeur du contenu
		\item \keyword{height} : la hauteur du contenu
		\item pour spécifier la direction des marges (internes ou externes) ou la bordure, on ajoute \keyword{-top} (haut), \keyword{-left} (gauche), \keyword{-right} (droit) ou \keyword{-bottom} (bas)
	\end{itemize}
\end{minipage}
%
\begin{minipage}{0.38\textwidth}
	\hgraphpage{box-model.pdf}
	\begin{exampleblock}{box.html}
		\lstset{escapeinside=**}
		\tiny\bfseries
\begin{lstlisting}[language={html}]
#d1 {
  background: red;
  width: 200px; height: 100px;
  padding: 5px;
  border-top: 5px solid green;
  border-bottom: 5px dashed green;
  margin-bottom: 20px;
}
#d1 p {background: blue;}
\end{lstlisting}
	\end{exampleblock}
\end{minipage}
\end{frame}

\subsection{Flotteurs et débordement}

\begin{frame}[fragile]
\frametitle{CSS : Mise en page}
\framesubtitle{Flotteurs et débordement : Flotteurs}

\begin{minipage}{0.60\textwidth}
	\begin{itemize}
		\item créer des mises en page de sites web avec plusieurs colonnes d'informations flottantes les unes à côté des autres
		\item éléments de type bloc
		\item l'attribut \keyword{float} avec la valeur \optword{left} (à gauche) ou \optword{right} (à droite)
		\item pour terminer le flottement, utiliser l'attribut \keyword{clear} avec la valeur \optword{both} (les deux types de flottement), \optword{left} (à gauche) ou \optword{right} (à droite)
	\end{itemize}
\end{minipage}
%
\begin{minipage}{0.38\textwidth}
	\begin{exampleblock}{float.html}
		\lstset{escapeinside=**}
		\scriptsize\bfseries
\begin{lstlisting}[language={html}]
.c1 {background: red;}
.c2 {background: blue;}
.c2::after {
  content: "";
  clear: both;
  display: block;
}
.c3 {background: green;}
#f {
  width: 200px;
  height: 200px;
  float: left;
  background: cyan;
}
\end{lstlisting}
	\end{exampleblock}
\end{minipage}
\end{frame}

\begin{frame}[fragile]
\frametitle{CSS : Mise en page}
\framesubtitle{Flotteurs et débordement : Débordement}

\begin{minipage}{0.60\textwidth}
	\begin{itemize}
		\item Dans un élément en bloc, si le contenu dépasse la taille de l'élément, on peut spécifier comment le gérer
		\item on utilise l'attribut \keyword{overflow} avec la valeur
		\begin{itemize}
			\item \optword{visible} le contenu n'est pas rogné (par défaut)
			\item \optword{hidden} le contenu est rogné
			\item \optword{scroll} ajouter des barres de défilement
			\item \optword{auto} comportement de l'agent utilisateur
		\end{itemize}
		\item On peut spécifier le dépassement vertical en utilisant l'attribut \keyword{overflow-y}
		\item On peut spécifier le dépassement horizontal en utilisant l'attribut \keyword{overflow-x}
	\end{itemize}
\end{minipage}
%
\begin{minipage}{0.38\textwidth}
	\begin{exampleblock}{overflow.html}
		\lstset{escapeinside=**}
		\tiny\bfseries
\begin{lstlisting}[language={html}]
div {
  width: 150px; height: 150px;
}

#d1 {
  background: yellow;
  overflow: scroll; }

#d2 {
  background: red;
  overflow-x: scroll;
  overflow-y: hidden; }

#d3 {
  background: blue;
  overflow: hidden; }

#d4 {
  background: orange;
  overflow: visible; }
\end{lstlisting}
	\end{exampleblock}
\end{minipage}
\end{frame}

\subsection{Position}

\begin{frame}[fragile]
\frametitle{CSS : Mise en page}
\framesubtitle{Position : Position des éléments (1)}

\begin{itemize}
	\item on utilise la propriété \keyword{position} avec la valeur
	\begin{itemize}
		\item \optword{static} L'élément est positionné dans le flux avec sa position. (par défaut)
		
		\item \optword{relative} L'élément est positionné dans le flux normal du document puis décalé, par rapport à lui-même.
		
		\item \optword{absolute} L'élément est retiré du flux normal et aucun espace n'est créé pour l'élément sur la page. Il est ensuite positionné par rapport à son ancêtre le plus proche qui a la position \optword{absolute} ou \optword{fixed}.
		
		\item \optword{fixed} L'élément est retiré du flux normal et aucun espace n'est laissé pour l'élément. L'élément est positionné relativement au bloc englobant initial formé par la zone d'affichage (viewport).
		
		\item \optword{sticky} La position de la boîte est calculée en fonction du flux normal du document. Ensuite, la boîte est décalée par rapport à son ancêtre de défilement le plus proche et par rapport à son bloc englobant. 
	\end{itemize}
\end{itemize}

\end{frame}

\begin{frame}[fragile]
\frametitle{CSS : Mise en page}
\framesubtitle{Position : Position des éléments (2)}

\begin{minipage}{0.60\textwidth}
	\begin{itemize}
		\item Les éléments avec position calculée : \optword{relative}, \optword{absolute}, \optword{fixed} et \optword{sticky}.
		\item La position est calculée selon les propriétés :
		\begin{itemize}
			\item \keyword{top} la distance en haut
			\item \keyword{left} la distance à gauche
			\item \keyword{right} la distance à droite
			\item \keyword{bottom} la distance en bas
		\end{itemize}
%		\item On a présenté le positionnement sur 2 dimensions (X et Y)
		\item Si on veut positionner une boîte sur une autre, on doit utiliser l'axe Z 
		\item L'attribut \keyword{z-index} a la valeur \optword{0} par défaut
		\item Une boîte est affichée sur une autre si son \keyword{z-index} est supérieur
	\end{itemize}
\end{minipage}
%
\begin{minipage}{0.38\textwidth}
	\begin{exampleblock}{position.html}
		\lstset{escapeinside=**}
		\tiny\bfseries\vspace{-6pt}
\begin{lstlisting}[language={CSS}]
.boite { width: 150px; 
  height: 150px;}
  
#stat { position: static;
  background: red;}
  
#rel { position: relative;
  background: green; z-index: 1;
  top: -50px; left:50px; }
  
#fix { position: fixed;
  top: 10px; right: 10px;
  background: yellow;
  height:200px; width:300px; }
  
#abs { position: absolute;
  background: blue;
  top: 2em; left:100px; }

#stic { position: sticky;
  background: magenta;
  top: 10px; left:300px; }
\end{lstlisting}\vspace{-6pt}
	\end{exampleblock}
\end{minipage}
\end{frame}


\begin{frame}[fragile]
\frametitle{CSS : Mise en page}
\framesubtitle{Position : Alignement pour une disposition flexible (flexbox)}

\begin{minipage}{0.60\textwidth}
	\begin{itemize}
		\item aligner horizontalement les bloques : \keyword{margin:auto;}
		\item espaces horizontalse entre les fils : \keyword{justify-content}
		\begin{itemize}
			\item \optword{flex-start} regrouper au début
			\item \optword{flex-end} regrouper à la fin
			\item \optword{center} regrouper au centre
			\item \optword{space-between} espacer équitablement
		\end{itemize}
		\item aligner verticalement les fils : \keyword{align-items}, soi-même : \keyword{align-self}
		\begin{itemize}
			\item \optword{flex-start} aligner en haut
			\item \optword{flex-end} aligner en bas
			\item \optword{center} centrer
			\item \optword{stretch} étirer
		\end{itemize}
	\end{itemize}
\end{minipage}
%
\begin{minipage}{0.38\textwidth}
	\begin{exampleblock}{alignement.html}
		\lstset{escapeinside=**}
		\tiny\bfseries\vspace{-6pt}
\begin{lstlisting}[language={html}]
.c {
  width: 300px;
  height: 200px;
  border: 2px dashed red;
  display: flex;
  align-items: flex-start;
  justify-content: space-between;
}

.c div {
  margin: 5px;
  background: yellow;
}

#deux {
  align-self: stretch;
}

#trois {
  margin-top: auto;
  margin-bottom: auto;
}
\end{lstlisting}\vspace{-6pt}
	\end{exampleblock}
\end{minipage}
\end{frame}

\subsection{Réactivité (responsiveness)}

\begin{frame}[fragile]
\frametitle{CSS : Mise en page}
\framesubtitle{Réactivité (responsiveness) : Taille de la fenêtre}

%\begin{minipage}{0.60\textwidth}
	\begin{itemize}
		\item les navigateurs mobiles ont tendance à mentir sur leurs largeurs de la fenêtre.
		\item pour optimiser la taille des mobiles, utiliser le méta \keyword{viewport} dans l'entête HTML
		\item \keyword{width=device-width} indique aux navigateurs mobiles qu'ils doivent définir la largeur de la fenêtre d'affichage à la largeur de l'appareil
		\item \keyword{initial-scale=1} : mettre le document à 100\% de sa taille prévue
	\end{itemize}
%\end{minipage}
%
%\begin{minipage}{0.38\textwidth}
	\begin{block}{Taille de  la fenêtre}
		\lstset{escapeinside=**}
		\scriptsize\bfseries
\begin{lstlisting}[language={html}]
<head>
  <meta name="viewport" content="width=device-width,initial-scale=1.0">
</head>
\end{lstlisting}
	\end{block}
%\end{minipage}
\end{frame}

\begin{frame}[fragile]
\frametitle{CSS : Mise en page}
\framesubtitle{Réactivité (responsiveness) : Requêtes média (media queries)}

\begin{minipage}{0.60\textwidth}
\begin{itemize}
	\item \keyword{@media} avec des règles dedans
	\item Les types des médias
	\begin{itemize}
		\item \keyword{all} tous les appareils
		\item \keyword{print} impression
		\item \keyword{screen} appareils avec écran
		\item \keyword{speech} outils de synthèse vocale
	\end{itemize}
	\item Les caractéristiques des médias
	\begin{itemize}
		\item \keyword{width} la largeur (\keyword{min-width}, \keyword{max-width})
		\item \keyword{height} la hauteur (\keyword{min-height}, \keyword{max-height})
		\item \keyword{orientation} l'orientation, deux valeurs : \optword{portrait} et \optword{landscape}
		\item \keyword{resolution} densité de pixel 
	\end{itemize}
\end{itemize}
\end{minipage}
%
\begin{minipage}{0.38\textwidth}
\begin{exampleblock}{responsive.html}
	\lstset{escapeinside=**}
	\tiny\bfseries\vspace{-6pt}
\begin{lstlisting}[language={CSS}]
.c {
  background: green;
  height: 100px; width: auto;
}

@media screen and (min-width:640px) {
  .c {
    background: blue;
    height: 200px; width: 560px;
    margin-left: auto;
    margin-right: auto;
  }
}

@media print {
  .c {
    background: none;
    height: 2cm;
    border: 2mm dashed red;
  }
}
\end{lstlisting}\vspace{-6pt}
\end{exampleblock}
\end{minipage}
\end{frame}

%===================================================================================
\section{Propriétés de texte}
%===================================================================================

\subsection{Texte}

\begin{frame}[fragile]
\frametitle{CSS : Propriétés de texte}
\framesubtitle{Texte : Mise en page du texte}

\begin{minipage}{0.60\textwidth}
	\begin{itemize}
		\item \keyword{text-align} : alignement dans un élément de type bloque
		\begin{itemize}
			\item \optword{left} : aligner à gauche
			\item \optword{right} : aliger à droite
			\item \optword{center} : centrer
			\item \optword{justify} : justifier 
		\end{itemize}
		\item \keyword{text-indent} : retrait de la première ligne
		\item \keyword{line-height} : hauteur minimale des lignes
		\item \keyword{direction} : direction du texte
		\begin{itemize}
			\item \optword{ltr} : de gauche à droite
			\item \optword{rtl} : de droite à gauche 
		\end{itemize}
	\end{itemize}
\end{minipage}
%
\begin{minipage}{0.38\textwidth}
	\begin{exampleblock}{texte.html}
		\lstset{escapeinside=**}
		\scriptsize\bfseries\vspace{-6pt}
\begin{lstlisting}[language={CSS}]
p {
  width: 400px; 
  padding: 5px;
}

p#p1 {
  background: yellow;
  text-align: justify;
  text-indent: 1.5cm;
  line-height: 2em;
}

p#p2 {
  background: cyan;
  direction: rtl;
}
\end{lstlisting}\vspace{-6pt}
	\end{exampleblock}
\end{minipage}
\end{frame}

\begin{frame}[fragile]
\frametitle{CSS : Propriétés de texte}
\framesubtitle{Texte : Orientation}

\begin{minipage}{0.60\textwidth}
	\begin{itemize}
		\item \keyword{writing-mode} : mode d'écriture
		\begin{itemize}
			\item \optword{horizontal-tb} : de gauche à droite puis de haut en bas
			\item \optword{vertical-rl} : de haut en bas puis de droite à gauche
			\item \optword{vertical-lr} : de haut en bas puis de gauche à droite
		\end{itemize}
%		\item \keyword{text-orientation} : orientation du texte en mode d'écriture vertical
%		\begin{itemize}
%			\item \optword{mixed} : tourner les caractères de 90\textdegree dans le sens des aiguilles d'une montre. 
%			Les scripts verticaux sont affichés normalement. (par défaut)
%			\item \optword{upright} : afficher les caractères un après l'autre verticalement « debout ». 
%			\item \optword{sideways} : tourner les caractères de 90\textdegree dans le sens des aiguilles d'une montre pour tous les scripts.
%		\end{itemize}
	\end{itemize}
\end{minipage}
%
\begin{minipage}{0.38\textwidth}
	\begin{exampleblock}{texte.html}
		\lstset{escapeinside=**}
\tiny\bfseries\vspace{-6pt}
\begin{lstlisting}[language={CSS}]
.ib {
  height: 150px;
  display: inline-block;
  border: 2px solid black;
  margin-right: 10px;
}
.vlr {writing-mode: vertical-lr;}
.vrl {writing-mode: vertical-rl;}
.htb {writing-mode: horizontal-tb;}
.mixed {text-orientation: mixed;}
.upright {text-orientation: upright;}
.sideways {text-orientation: sideways;}
\end{lstlisting}\vspace{-6pt}
	\end{exampleblock}
\end{minipage}
\begin{itemize}
\item \keyword{text-orientation} : orientation du texte en mode d'écriture vertical
\begin{itemize}
	\item \optword{mixed} : tourner les caractères de 90\textdegree dans le sens des aiguilles d'une montre. 
	Les scripts verticaux sont affichés normalement. (par défaut)
	\item \optword{upright} : afficher les caractères un après l'autre verticalement « debout ». 
	\item \optword{sideways} : tourner les caractères de 90\textdegree dans le sens des aiguilles d'une montre pour tous les scripts.
\end{itemize}
\end{itemize}
\end{frame}

\begin{frame}[fragile]
\frametitle{CSS : Propriétés de texte}
\framesubtitle{Texte : Casse et décoration}

\begin{minipage}{0.60\textwidth}
	\begin{itemize}
		\item \keyword{text-transform} : définir la casse du texte
		\begin{itemize}
			\item \optword{uppercase} : tous en majuscule
			\item \optword{lowercase} : tous en minuscule
			\item \optword{capitalize} : les premières lettres en majuscule
		\end{itemize}
		\item \keyword{text-decoration} :  les décorations de texte
		\begin{itemize}
			\item \optword{underline} : souligner le texte
			\item \optword{overline} : tracer une ligne sur le texte
			\item \optword{line-through} : barrer le texte
			\item type de ligne
			\item couleur de ligne
		\end{itemize}
		\item \keyword{text-shadow} : ombres du texte
		\begin{itemize}
			\item \expword{décalage-x décalage-y rayon couleur} 
		\end{itemize}
	\end{itemize}
\end{minipage}
%
\begin{minipage}{0.38\textwidth}
	\begin{exampleblock}{texte.html}
		\lstset{escapeinside=**}
		\tiny\bfseries
\begin{lstlisting}[language={CSS}]
.maj {
  text-transform: uppercase;
}

.min {
  text-transform: lowercase;
}

.cap {
  text-transform: capitalize;
}

.dec1 {
  text-decoration: underline wavy red;
  text-shadow: 2px 2px 3px yellow;
}

.dec2 {
  text-decoration: underline overline;
}
\end{lstlisting}
	\end{exampleblock}
\end{minipage}
\end{frame}

\subsection{Police}

\begin{frame}[fragile]
\frametitle{CSS : Propriétés de texte}
\framesubtitle{Police}

\begin{minipage}{0.60\textwidth}
	\begin{itemize}
		\item \keyword{font-family} : la police utilisée
		\item \keyword{font-size} : la taille de la police
		\item \keyword{font-style} : le style de la police
		\item \keyword{font-weight} : 
	\end{itemize}
\end{minipage}
%
\begin{minipage}{0.38\textwidth}
	\begin{exampleblock}{Structure d'un document HTML}
		\lstset{escapeinside=**}
		\scriptsize\bfseries
		\begin{lstlisting}[language={html}]
		h1 {
		color: red;
		}
		\end{lstlisting}
	\end{exampleblock}
\end{minipage}
\end{frame}

\begin{frame}[fragile]
\frametitle{CSS : Propriétés de texte}
\framesubtitle{Police : Définir une famille de police}

\begin{minipage}{0.60\textwidth}
	\begin{itemize}
		\item \keyword{@font-face}
	\end{itemize}
\end{minipage}
%
\begin{minipage}{0.38\textwidth}
\begin{exampleblock}{Structure d'un document HTML}
\lstset{escapeinside=**}
\scriptsize\bfseries
\begin{lstlisting}[language={html}]
@font-face {
  font-family: "maPolice";
  src: url("unePolice.ttf");
}
\end{lstlisting}
\end{exampleblock}
\end{minipage}
\end{frame}

\subsection{Mots et lettres}

\begin{frame}[fragile]
\frametitle{CSS : Propriétés de texte}
\framesubtitle{Mots et lettres}

\begin{minipage}{0.60\textwidth}
	\begin{itemize}
		\item \keyword{word-spacing}
		\item \keyword{word-break}
		\item \keyword{letter-spacing}
	\end{itemize}
\end{minipage}
%
\begin{minipage}{0.38\textwidth}
	\begin{exampleblock}{Structure d'un document HTML}
		\lstset{escapeinside=**}
		\scriptsize\bfseries
		\begin{lstlisting}[language={html}]
		h1 {
		color: red;
		}
		\end{lstlisting}
	\end{exampleblock}
\end{minipage}
\end{frame}

\subsection{Listes}

\begin{frame}[fragile]
\frametitle{CSS : Propriétés de texte}
\framesubtitle{Listes}

\begin{minipage}{0.60\textwidth}
	\begin{itemize}
		\item \keyword{list-style-type}
		\item \keyword{list-style-position}
		\item \keyword{list-style-image}
	\end{itemize}
\end{minipage}
%
\begin{minipage}{0.38\textwidth}
	\begin{exampleblock}{Structure d'un document HTML}
		\lstset{escapeinside=**}
		\scriptsize\bfseries
		\begin{lstlisting}[language={html}]
		h1 {
		color: red;
		}
		\end{lstlisting}
	\end{exampleblock}
\end{minipage}
\end{frame}

%===================================================================================
\section{Aspect et animations}
%===================================================================================

\subsection{Arrière-plan}

\begin{frame}[fragile]
\frametitle{CSS : Aspect et animations}
\framesubtitle{Arrière-plan}

\begin{minipage}{0.60\textwidth}
	\begin{itemize}
		\item \keyword{background-color}
		\item \keyword{background-image}
		\item \keyword{background-repeat}
		\item \keyword{background-attachment}
	\end{itemize}
\end{minipage}
%
\begin{minipage}{0.38\textwidth}
	\begin{exampleblock}{Structure d'un document HTML}
		\lstset{escapeinside=**}
		\scriptsize\bfseries
		\begin{lstlisting}[language={html}]
		h1 {
		color: red;
		}
		\end{lstlisting}
	\end{exampleblock}
\end{minipage}
\end{frame}

\begin{frame}[fragile]
\frametitle{CSS : Aspect et animations}
\framesubtitle{Arrière-plan : Arrière-plan dégradé}

\begin{minipage}{0.60\textwidth}
	\begin{itemize}
		\item 
	\end{itemize}
\end{minipage}
%
\begin{minipage}{0.38\textwidth}
	\begin{exampleblock}{Structure d'un document HTML}
		\lstset{escapeinside=**}
		\scriptsize\bfseries
		\begin{lstlisting}[language={html}]
		h1 {
		color: red;
		}
		\end{lstlisting}
	\end{exampleblock}
\end{minipage}
\end{frame}

\subsection{Transformations}

\begin{frame}[fragile]
\frametitle{CSS : Aspect et animations}
\framesubtitle{Transformations}

\begin{minipage}{0.60\textwidth}
	\begin{itemize}
		\item 
	\end{itemize}
\end{minipage}
%
\begin{minipage}{0.38\textwidth}
	\begin{exampleblock}{Structure d'un document HTML}
		\lstset{escapeinside=**}
		\scriptsize\bfseries
		\begin{lstlisting}[language={html}]
		h1 {
		color: red;
		}
		\end{lstlisting}
	\end{exampleblock}
\end{minipage}
\end{frame}

\subsection{Animations}

\begin{frame}[fragile]
\frametitle{CSS : Aspect et animations}
\framesubtitle{Animations}

\begin{minipage}{0.60\textwidth}
	\begin{itemize}
		\item 
	\end{itemize}
\end{minipage}
%
\begin{minipage}{0.38\textwidth}
	\begin{exampleblock}{Structure d'un document HTML}
		\lstset{escapeinside=**}
		\scriptsize\bfseries
		\begin{lstlisting}[language={html}]
		h1 {
		color: red;
		}
		\end{lstlisting}
	\end{exampleblock}
\end{minipage}
\end{frame}


\begin{frame}
\frametitle{CSS}
\framesubtitle{Un peu d'humour}

\begin{center}
	\vgraphpage{css-humour.jpg}
\end{center}

\end{frame}

\insertbibliography{Bweb07}{*}

\end{document}

